
\documentclass{article}
\usepackage{colortbl}
\usepackage{makecell}
\usepackage{multirow}
\usepackage{supertabular}

\begin{document}

\newcounter{utterance}

\centering \large Interaction Transcript for game `cladder', experiment `full\_v1.5\_default', episode 7019 with Qwen3.5{-}4B.
\vspace{24pt}

{ \footnotesize  \setcounter{utterance}{1}
\setlength{\tabcolsep}{0pt}
\begin{supertabular}{c@{$\;$}|p{.15\linewidth}@{}p{.15\linewidth}p{.15\linewidth}p{.15\linewidth}p{.15\linewidth}p{.15\linewidth}}
    \# & \multicolumn{2}{c}{Player} && \multicolumn{2}{c}{Game Master} \\
    \hline

    \theutterance \stepcounter{utterance}  
    & & & \multicolumn{4}{p{0.6\linewidth}}{
        \cellcolor[rgb]{0.9,0.9,0.9}{
            \makecell[{{p{\linewidth}}}]{
                \texttt{\tiny{[P1$\langle$GM]}}
                \texttt{You are an expert in causal inference. The following question is not a typical commonsense query, but rather a meticulously designed question created by a professor specializing in causal inference, intended to assess the students' mastery of the course content.} \\
\texttt{Imagine a self{-}contained, hypothetical world with only the following conditions, and without any unmentioned factors or causal relationships: Hwax has a direct effect on gyzp and jyka. Jyka has a direct effect on lirg. Gyzp has a direct effect on lirg. Method 1: We look at how jyka correlates with lirg case by case according to gyzp. Method 2: We look directly at how jyka correlates with lirg in general. To understand how jyka affects lirg, is it more correct to use the Method 1 than Method 2?} \\
\texttt{Start your answer with 'Yes' or 'No', followed by additional reasoning or evidence to support your explanation.} \\
            }
        }
    }
     \\ \\

    \theutterance \stepcounter{utterance}  
    & \multicolumn{4}{p{0.6\linewidth}}{
        \cellcolor[rgb]{0.561,0.737,0.561}{
            \makecell[{{p{\linewidth}}}]{
                \texttt{\tiny{[P1$\rangle$GM]}}
                \texttt{Yes} \\
\\ 
\texttt{**Reasoning:**} \\
\\ 
\texttt{To determine which method is more correct for understanding the causal effect of *jyka* on *lirg*, we must analyze the causal structure and the presence of confounding variables.} \\
\\ 
\texttt{1.  **Causal Structure Analysis**:} \\
\texttt{    *   **Hwax** \$\textbackslash{}rightarrow\$ **Gyzp**} \\
\texttt{    *   **Hwax** \$\textbackslash{}rightarrow\$ **Jyka**} \\
\texttt{    *   **Jyka** \$\textbackslash{}rightarrow\$ **Lirg**} \\
\texttt{    *   **Gyzp** \$\textbackslash{}rightarrow\$ **Lirg**} \\
\\ 
\texttt{    In this diagram, *Hwax* is a common cause (a confounder) for both *jyka* and *lirg*. *Gyzp* is a collider or a mediator depending on the path, but crucially, *Gyzp* is a direct cause of *lirg* and shares a common cause (*Hwax*) with *jyka*.} \\
\\ 
\texttt{2.  **Analysis of Method 2 (General Correlation)**:} \\
\texttt{    Method 2 looks at the raw correlation between *jyka* and *lirg*. Because both variables are influenced by *Hwax*, they are statistically associated even if *jyka* had no causal effect on *lirg*. This is the definition of **confounding**. The observed correlation in Method 2 is a mixture of the true causal effect (\$Jyka \textbackslash{}rightarrow Lirg\$) and the spurious correlation induced by *Hwax*. Therefore, Method 2 yields a biased estimate of the causal effect.} \\
\\ 
\texttt{3.  **Analysis of Method 1 (Stratification by Gyzp)**:} \\
\texttt{    Method 1 examines the correlation between *jyka* and *lirg* within strata (subgroups) defined by the value of *gyzp*.} \\
\texttt{    *   While *gyzp* is a descendant of *hwx*, conditioning on a variable that is a descendant of a confounder can sometimes introduce bias (collider bias) if the conditioning is on a collider. However, in this specific chain (\$Hwax \textbackslash{}rightarrow Gyzp \textbackslash{}rightarrow Lirg\$), *Gyzp* is not a collider; it is a mediator on the path from *Hwax* to *Lirg*.} \\
\texttt{    *   More importantly, the question asks which method is *more* correct. Method 1 attempts to isolate the relationship by holding *gyzp* constant. By fixing *gyzp*, we remove the variation in *lirg* that is driven by *gyzp*.} \\
\texttt{    *   However, there is a nuance: If we condition on *gyzp*, we are essentially looking at the effect of *jyka* on *lirg* *given* a specific level of *gyzp*. Since *gyzp* is a direct cause of *} \\
            }
        }
    }
    & & \\ \\

    \theutterance \stepcounter{utterance}  
    & & & \multicolumn{2}{p{0.3\linewidth}}{
        \cellcolor[rgb]{0.9,0.9,0.9}{
            \makecell[{{p{\linewidth}}}]{
                \texttt{\tiny{[GM$|$GM]}}
                \texttt{Yes} \\
            }
        }
    }
    & & \\ \\

    \theutterance \stepcounter{utterance}  
    & & & \multicolumn{2}{p{0.3\linewidth}}{
        \cellcolor[rgb]{0.9,0.9,0.9}{
            \makecell[{{p{\linewidth}}}]{
                \texttt{\tiny{[GM$|$GM]}}
                \texttt{yes} \\
            }
        }
    }
    & & \\ \\

    \theutterance \stepcounter{utterance}  
    & & & \multicolumn{2}{p{0.3\linewidth}}{
        \cellcolor[rgb]{0.9,0.9,0.9}{
            \makecell[{{p{\linewidth}}}]{
                \texttt{\tiny{[GM$|$GM]}}
                \texttt{game\_result = WIN} \\
            }
        }
    }
    & & \\ \\

\end{supertabular}
}

\end{document}
